\chapter{随机事件与概率}\label{ch:01}
\chaptermap{样本空间\\事件运算}{等可能与几何\\公理化概率}{条件与分层\\全概率Bayes}{独立试验\\计数与递推}
\begin{examtip}
\textbf{考研核心：}事件表达、加法公式、条件概率、全概率、Bayes与独立性。
常考信号是“至少一次”“已知发生”“来自哪一类”。\textbf{竞赛视角：}容斥、对称、状态递推。
\textbf{测度关联：}事件必须属于可测集族，连续模型中概率为零不等于不可能。
\end{examtip}

\section{事件运算与概率公理}\label{sec:11}
\src{茆诗松, §1.1—1.3；李东风, §1.1、§1.3}
\subsection{概念与结论}
\begin{definition}{事件与概率空间}
样本空间$\Omega$是非空的可能结果集合。事件域$\Fcal$包含$\Omega$，对补集与可数并封闭，称为$\sigma$代数；事件是$\Fcal$中的集合。
$A\cup B$表示至少一个发生，$A\cap B$表示同时发生，$A\setminus B$表示$A$发生而$B$不发生。
概率是映射$\Pp:\Fcal\to[0,1]$，满足$\Pp(\Omega)=1$，且对两两不交的$A_n\in\Fcal$有
$\Pp(\bigcup_{n\ge1}A_n)=\sum_{n\ge1}\Pp(A_n)$。三者组成概率空间$(\Omega,\Fcal,\Pp)$。
\end{definition}
\begin{theorem}{加法、连续性与并集界}
$\Pp(A^c)=1-\Pp(A)$，$\Pp(A\cup B)=\Pp(A)+\Pp(B)-\Pp(A\cap B)$。
若$A_n\uparrow A$或$A_n\downarrow A$，则$\Pp(A_n)\to\Pp(A)$。
一般有$\Pp(\bigcup_n A_n)\le\sum_n\Pp(A_n)$。
\end{theorem}
\begin{proof}
把$A\cup B$拆成$A\setminus B$、$A\cap B$、$B\setminus A$即得加法公式。
对递增事件写$A=A_1\sqcup\bigsqcup_{n\ge2}(A_n\setminus A_{n-1})$，可数可加后取部分和得连续性；递减情形取补集。
任意并集先按“首次出现的下标”拆成不交事件，再逐项放大为$A_n$，即得并集界。
\end{proof}
\begin{method}{把语言翻译成互斥事件}
适用：多个事件的“仅”“至少”“恰好”。先写集合式，再用Venn区域分解，最后计算概率。
“仅$A$发生”为$A\cap B^c\cap C^c$；“恰好一个”是三个此类区域之并。
易错点：互斥是集合关系，独立是概率关系。
\end{method}
\subsection{分层例题}
\worked{恰好一个事件}{基础例题}{\star}{1-1}
{设$\Pp(A)=0.6,\Pp(B)=0.5,\Pp(A\cap B)=0.2$，求恰好一个发生的概率。}
{找对称差的两个不交区域。}
{方法一：$\Pp(A\setminus B)+\Pp(B\setminus A)=(0.6-0.2)+(0.5-0.2)=0.7$。
方法二：先算至少一个的概率$0.6+0.5-0.2=0.9$，再减同时发生的$0.2$，仍得$0.7$。}
{分区域法便于检查；先并后减适合已有并集概率时使用。}
{直接把$0.6+0.5$当答案会重复计入交集两次。}
\worked{三事件容斥}{考研风格}{\star\star}{1-2}
{三事件各有概率$\frac{1}{2}$，两两交概率均为$\frac{1}{4}$，三者交概率$\frac{1}{8}$，求至少一个发生的概率。}
{逐级纠正重复计数。}
{$\Pp(A\cup B\cup C)=\frac{3}{2}-\frac{3}{4}+\frac{1}{8}=\frac{7}{8}$。交中三者同时发生的区域在第一轮计3次、第二轮减3次，最后必须补1次。
本题数据也满足相互独立，补事件概率为$\frac{1}{8}$可作检验。}
{容斥只依赖各阶交概率，不预先要求独立。}
{只知道两两独立时，不能自动确定三重交概率。}
\worked{无放回排列的固定点}{竞赛风格}{\star\star\star}{1-3}
{随机排列$1,2,3,4$，求至少一数留在原位的概率。}
{令$A_i$为第$i$个数留在原位，对固定点使用容斥。}
{固定指定$k$个位置有$(4-k)!$种排列，故所求为
\[\dfrac{\binom41 3!-\binom42 2!+\binom43 1!-\binom44 0!}{4!}=\dfrac{15}{24}=\dfrac58.\]
其补集错排数为$9$，总数$24$。每次容斥按指定位置集合计数，因此多固定点排列的重复次数得到正确消除。}
{把多个“至少”转成容斥，复杂事件变成容易计算的交事件。}
{推广至$n$时，无固定点概率为$\sum_{k=0}^n\frac{(-1)^k}{k!}$。}

\section{古典概型与几何概型}\label{sec:12}
\src{茆诗松, §1.2；李东风, §1.2}
\begin{definition}{等可能与均匀模型}
有限等可能空间中$\Pp(A)=\frac{|A|}{|\Omega|}$。在有限正体积区域$D\subset\R^d$上均匀取点时，$\Pp(A)=\frac{\operatorname{vol}(A\cap D)}{\operatorname{vol}(D)}$。
均匀指相对于指定坐标下的体积；更换随机生成机制可能改变概率。
\end{definition}
\begin{method}{先统一计数口径}
适用：抽球、排列、抽签与几何随机点。先定有序或无序、放回或不放回；分子分母保持同一口径。
几何题先画可行域，再画事件域；均匀性未经说明时不能直接用面积比。
\end{method}
\worked{两次抽到红球}{基础例题}{\star}{1-4}
{盒中3红2白，不放回抽2球，求两球都红的概率。}
{用无序组合与逐步条件概率互证。}
{方法一：$\frac{\binom32}{\binom52}=\frac{3}{10}$。
方法二：按抽取顺序，$\Pp(R_1R_2)=\frac35\frac24=\frac{3}{10}$。
前者把一个无序球对作为样本点，后者把有序球对作为样本点，总数分别为10和20。}
{组合式短；乘法式便于处理逐次条件变化。}
{已知第一球为红时，第二球为红的条件概率为$\frac{2}{4}$；第二球为红的无条件概率仍为$\frac{3}{5}$。}
\worked{相遇问题}{考研风格}{\star\star}{1-5}
{两人独立在一小时内均匀到达，各等待15分钟，求相遇概率。}
{以小时为单位，研究单位正方形内$|x-y|\le\frac{1}{4}$。}
{不相遇区域是两块直角边长$\frac{3}{4}$的直角三角形，故概率
$1-2\cdot\frac12(\frac{3}{4})^2=\frac{7}{16}$。
另一法按第一人到达时间$x$条件化。允许的第二人到达时间区间为$[x-\frac{1}{4},x+\frac{1}{4}]\cap[0,1]$，其长度分段积分为
\[\int_0^{\frac{1}{4}}(x+\dfrac{1}{4})\dd x+\int_{\frac{1}{4}}^{\frac{3}{4}}\dfrac12\dd x+
\int_{\frac{3}{4}}^1(\dfrac{5}{4}-x)\dd x=\dfrac3{32}+\dfrac14+\dfrac3{32}=\dfrac7{16}.\]}
{补域通常比相遇带状区域更容易算面积。}
{“各等待15分钟”对应时间差不超过15分钟，而非30分钟。}
\worked{三骰点数和}{考研风格}{\star\star}{1-6}
{三枚公平骰子独立投掷，求点数和为12的概率。}
{有序结果等可能；隔板法必须加入上界。}
{令$y_i=x_i-1$，则$y_1+y_2+y_3=9$，且$0\le y_i\le5$。
无上界解数为$\binom{11}{2}=55$；指定一个$y_i\ge6$时有$\binom52=10$种，两项同时超界不可能。因此有效数为$55-3\cdot10=25$，所求$\frac{25}{216}$。}
{容斥修正隔板法，把上下界都保留在模型中。}
{55只是无上界正整数解数，不能直接作骰子结果数。}
\begin{figure}[htbp]\centering
\begin{tikzpicture}[scale=4]
\fill[Teal!12](0,0)--(.25,0)--(1,.75)--(1,1)--(.75,1)--(0,.25)--cycle;
\draw[Ink](0,0) rectangle(1,1);
\draw[Gold](0,.25)--(.75,1);\draw[Gold](.25,0)--(1,.75);
\node[below]at(.5,0){$x$};\node[left]at(0,.5){$y$};\node at(.5,.5){$|x-y|\le\frac{1}{4}$};
\end{tikzpicture}\caption{相遇事件是正方形中的对角带}\label{fig:meeting}\end{figure}

\section{条件概率与乘法公式}\label{sec:13}
\src{茆诗松, §1.4；李东风, §1.4}
\begin{definition}{条件概率}
若$\Pp(B)>0$，定义$\Pp(A\given B)=\frac{\Pp(A\cap B)}{\Pp(B)}$。
于是$\Pp(A\cap B)=\Pp(B)\Pp(A\given B)$；多阶段反复使用即可得到链式乘法公式。
\end{definition}
\begin{method}{条件化先缩小分母}
适用：“已知至少”“观察到某结果”。先写给定事件$B$，把分母改为$\Pp(B)$；再算$A\cap B$。
若题目包含观察机制，必须把该机制加入事件，不可只按口头条件猜比例。
\end{method}
\worked{已知至少一次成功}{基础例题}{\star\star}{1-7}
{独立投两枚公平硬币，已知至少一正，求两枚都正的条件概率。}
{给定事件排除反反。}
{方法一：分母$1-\frac{1}{4}=\frac{3}{4}$，分子$\frac{1}{4}$，故答案$\frac{1}{3}$。
方法二：条件样本空间为正正、正反、反正，三点仍等可能，只有一点满足两正。
若条件改为“第一枚为正”，空间只有正正、正反，答案变为$\frac{1}{2}$。}
{分式法适用任意权重；缩空间计数要求剩余点仍等可能。}
{至少一正与指定一枚为正是不同条件。}
\worked{逐次抽球的条件}{考研风格}{\star\star}{1-8}
{3红2白不放回抽2球，已知第二球红，求第一球红的概率。}
{先求联合概率，再用抽取位置的对称性算边缘。}
{$\Pp(R_1R_2)=\frac{3}{10}$，$\Pp(R_2)=\frac{3}{5}$，故$\Pp(R_1\given R_2)=\frac{1}{2}$。
也可将第二个红球固定，从剩余4球中看第一球，剩下2红2白，比例为$\frac{1}{2}$。}
{逆时间条件也可按定义求解，不要求按观察顺序推理。}
{不能把$\Pp(R_1\given R_2)$与$\Pp(R_2)$混为一谈。}
\worked{无记忆性的离散刻画}{竞赛风格}{\star\star\star}{1-9}
{正整数值$N$满足$\Pp(N>m+n\given N>m)=\Pp(N>n)$且尾概率为正，证明$N$为几何分布。}
{把条件概率变成尾序列的乘法方程。}
{记$q_n=\Pp(N>n)$，$q_0=1$。条件给$q_{m+n}=q_mq_n$，所以$q_n=q_1^n$。
由于$N$有限几乎必然，$q_n\downarrow0$，故$0\le q_1<1$；严格正尾假设排除$q_1=0$。
令$p=1-q_1$，则$\Pp(N=k)=q_{k-1}-q_k=p(1-p)^{k-1}$。}
{把概率性质转成函数方程，是竞赛中常见的分布刻画。}
{取值从1开始；若记录首次成功前失败数，指数应整体平移。}

\section{全概率与Bayes公式}\label{sec:14}
\src{茆诗松, §1.4；李东风, §1.6}
\begin{theorem}{分层求和与反向更新}
设$B_i$构成可数划分且$\Pp(B_i)>0$，则
\begin{align}
\Pp(A)&=\sum_i\Pp(B_i)\Pp(A\given B_i),\\
\Pp(B_j\given A)&=\dfrac{\Pp(B_j)\Pp(A\given B_j)}{\sum_i\Pp(B_i)\Pp(A\given B_i)},\quad\Pp(A)>0.
\end{align}
\end{theorem}
\begin{proof}
$A=\bigsqcup_i(A\cap B_i)$，先可数求和再用乘法公式得第一式。
第二式按条件概率定义，把分母替换为第一式即可。
\end{proof}
\begin{method}{树图的正向与反向}
适用：机器来源、盒子选择、群体筛查。先验写在第一层，条件概率写在第二层；路径相乘，终点相加。
Bayes反查某路径时用该路径权重除以所有通向同一观察结果的权重。
\end{method}
\worked{混合生产线}{考研风格}{\star\star}{1-10}
{甲乙产量比$3:2$，次品率分别$0.01,0.03$。随机一件为次品，求来自乙的概率。}
{先求次品总概率。}
{方法一：$\Pp(D)=0.6\cdot0.01+0.4\cdot0.03=0.018$，后验为$\frac{0.012}{0.018}=\frac{2}{3}$。
方法二：设生产10000件，甲次品期望60件、乙120件，权重比例$60:120$，归一化后乙占$\frac{2}{3}$。
方法三：后验赔率$\frac{\Pp(\text{乙}\given D)}{\Pp(\text{甲}\given D)}=\frac{0.4}{0.6}\frac{0.03}{0.01}=2$。}
{公式通用；虚拟频数直观；赔率法适合两类反复更新。}
{后验既取决于次品率，也取决于产量先验。}
\worked{低基率筛查模型}{考研风格}{\star\star}{1-11}
{某类对象比例$0.01$，识别为阳性的概率为$0.9$；其他对象误阳性概率$0.05$。求阳性后属于该类的概率。}
{把真阳与误阳都计入分母。}
{真阳概率$0.01\cdot0.9=0.009$，误阳概率$0.99\cdot0.05=0.0495$，故后验$\frac{0.009}{0.0585}=\frac{2}{13}$。
若仅把$0.9$当后验，就忽视了数量更多的另一类。这里是抽象概率模型，不涉及实际检测数据。}
{Bayes训练中，先验小会显著影响观察后的判断。}
{条件方向不可倒置：$\Pp(+\given C)$不是$\Pp(C\given +)$。}
\worked{连续条件的初等对应}{测度论选学}{\star\star}{1-12}
{在单位正方形均匀取$(X,Y)$，求$\Pp(Y<X\given X>\frac{1}{2})$。}
{条件集合面积为$\frac{1}{2}$，在其内部积事件面积。}
{分子$\int_{\frac{1}{2}}^1x\dd x=\frac{3}{8}$，所以条件概率$\frac{3}{4}$。
也可在条件后先让$X$在$(\frac{1}{2},1)$均匀，给定$X=x$时$\Pp(Y<x)=x$，再平均得到$2\int_{\frac{1}{2}}^1x\dd x=\frac{3}{4}$。
后一种写法预示了条件期望的分层积分。}
{面积比与条件积分完全一致。}
{若条件换成$X=x$，分母概率为零，不能直接用事件条件概率分式。}

\section{独立性与Bernoulli试验}\label{sec:15}
\src{茆诗松, §1.5；Papoulis--Pillai, §3.2}
\begin{definition}{相互独立与Bernoulli概型}
事件族相互独立要求每个有限子族的交概率等于边缘概率之积。只核对两两相交不够。
$n$次独立同成功率$p$的试验中，成功数$N$满足$\Pp(N=k)=\binom nkp^k(1-p)^{n-k}$。
\end{definition}
\begin{proof}
指定$k$个位置成功的概率为$p^k(1-p)^{n-k}$，不同位置集合对应互斥事件，共$\binom nk$个，因此相加得二项分布。
\end{proof}
\begin{method}{至少、恰好与等待}
适用：独立重复试验。恰好$k$次用二项计数；至少一次优先补事件；连续模式或停止规则用状态递推。
先检查每次成功率是否相同以及试验是否独立，不能把无放回抽样直接当二项。
\end{method}
\worked{至少一次命中}{基础例题}{\star}{1-13}
{独立射击3次，每次命中率$\frac{1}{3}$，求至少一次命中的概率。}
{比较补事件、直接分类与递推。}
{方法一：$1-(\frac{2}{3})^3=\frac{19}{27}$。
方法二：$\sum_{k=1}^3\binom3k(\frac{1}{3})^k(\frac{2}{3})^{3-k}=\frac{12+6+1}{27}$。
方法三：记$a_n$为至少一次命中概率，$a_0=0$，按首轮分解$a_n=\frac{1}{3}+(\frac{2}{3})a_{n-1}$，推出$a_n=1-(\frac{2}{3})^n$。}
{单次问题补事件最短；递推适合推广变成功率或状态问题。}
{对立事件是“三次全失”，不是“恰有一次失误”。}
\worked{两两独立不等于相互独立}{竞赛风格}{\star\star}{1-14}
{独立投两枚公平硬币，$A$为第一枚正，$B$为第二枚正，$C$为两枚相同。验证两两独立而非相互独立。}
{逐个比较交概率。}
{三事件概率均为$\frac{1}{2}$；每个两两交只有正正或相应单点，概率$\frac{1}{4}$，所以两两独立。
但$A\cap B\cap C=\{\text{正正}\}$，其概率$\frac{1}{4}\ne(\frac{1}{2})^3$，故不相互独立。}
{反例通过第三事件编码前两次结果的关系。}
{三个两两等式只是相互独立所需条件的一部分。}
\worked{连续两正面的等待期望}{竞赛风格}{\star\star\star}{1-15}
{独立公平硬币投至首次连续两正，求投掷次数期望。}
{以当前连续正面的长度0或1为状态。}
{设剩余期望为$e_0,e_1$。从状态0投一次后，$e_0=1+\frac12e_1+\frac12e_0$；从状态1，正则结束、反则回0，所以$e_1=1+\frac12e_0$。
联立得$e_0=6,e_1=4$。期望有限可由每两次一组的“该组两正”成功率$\frac{1}{4}$控制等待时间，故方程求出的有限值有依据。}
{状态递推避免枚举所有停止序列。}
{等待两次成功与等待连续两次成功不是同一问题。}
\begin{warning}
互斥且概率都正的事件不独立；独立事件也未必互斥。
几何概型取决于抽样机制。连续空间的单点概率可以为零，仍是可能样本点。
Bayes分母必须汇总所有来源，不可只累加“有利”分支。
\end{warning}

\section*{图解、理解与巧解复盘}
\begin{insight}{概率空间先定权重，后做计数}
“随机”只说明结果不确定，不自动给出等可能。两枚骰子之和2与7分别有1与6个有序逆像，因此和值不能作为等权样本点。
三骰和为12的25个有序结果见例\ref{ex:1-6}；每个样本点权重均为$\frac{1}{216}$，这才是计数可用的原因。
\end{insight}
\begin{clevertrick}{至少一次的递推可转成乘法}
若每轮独立且成功率为$p_i$，记前$i$轮至少成功一次概率为$a_i$。全概率给$a_i=a_{i-1}+(1-a_{i-1})p_i$，所以$1-a_i=(1-a_{i-1})(1-p_i)$。
由$a_0=0$，立得$a_n=1-\prod_{i=1}^n(1-p_i)$。递推与补事件在“尚未成功的质量”上完全一致。
\end{clevertrick}
\begin{figure}[htbp]\centering
\begin{tikzpicture}[>=Stealth,every node/.style={font=\small},level distance=27mm]
\node(a)at(0,0){来源};\node(b)at(3,1){甲：$0.6$};\node(c)at(3,-1){乙：$0.4$};
\node(d)at(7,1){甲且次品：$0.006$};\node(e)at(7,-1){乙且次品：$0.012$};
\draw[->,Teal,thick](a)--(b);\draw[->,Teal,thick](a)--(c);
\draw[->,Gold,thick](b)--node[above]{$0.01$}(d);\draw[->,Gold,thick](c)--node[below]{$0.03$}(e);
\draw[decorate,decoration={brace,amplitude=5pt},Ink](9,1.3)--(9,-1.3)node[midway,right=5mm]{$\text{乙:甲}=2:1$};
\end{tikzpicture}\caption{Bayes反查按路径联合权重归一化，乙的后验为$\frac{2}{3}$}\label{fig:bayes-tree}\end{figure}

\section*{章末分层习题}
\exercise{1-a}{基础}{已知$\Pp(A)=0.4,\Pp(B)=0.7,\Pp(A\cap B)=0.2$，求$\Pp(A^c\cap B)$和$\Pp(A\cup B)$。}
\exercise{1-b}{基础}{6张标号签不放回抽2张，求两张都是偶数的概率。}
\exercise{1-c}{考研风格}{两厂供货比例$1:3$，合格率分别$0.8,0.9$。已知产品不合格，求来自第一厂的概率。}
\exercise{1-d}{考研风格}{两人独立在$[0,1]$均匀到达，各等$\frac{1}{3}$时间，求相遇概率。}
\exercise{1-e}{竞赛拓展}{三次独立试验成功率$p$，求恰好没有相邻成功的概率。}
\exercise{1-f}{竞赛拓展}{公平硬币投至首次出现“正反”，用状态方程求期望次数。}
\exercise{1-g}{测度论选学}{在$[0,1]$均匀取点，证明有理数事件概率为零，说明概率零与空集的区别。}
\exercise{1-h}{测度论选学}{证明若$\Pp(A_n)=0$对所有$n$成立，则$\Pp(\bigcup_n A_n)=0$；说明为何“可数”不可随意改成“任意”。}
\section*{章末方法速查}
\begin{quicktable}
容斥与补事件&已知交概率或反面简单&至少、恰好&遗漏高阶交项\\
计数与面积&等可能或均匀已确认&抽取、随机点&分子分母口径不同\\
全概率与Bayes&可测划分、分母正&来源、已知结果&反向条件与先验遗漏\\
状态递推&未来只依赖当前状态&连续模式、等待&状态不完整或期望未证有限\\
\end{quicktable}
